\section{Implementation and Evaluation Details}
\label{app:implementation}

\subsection{Inputs, checkpoints, and reproducibility}
The evaluation unit is a task--input pair, not necessarily a distinct source molecule across tasks. The common manifest comprises IND tasks BDP, BDQ, BPQ, DPQ, and BDPQ (500 pairs each), and OOD tasks MPQ (500), BDMQ (500), BHMQ (118), BMPQ (191), and HMPQ (96). Thus every method--model configuration has 3,905 evaluated outputs. B, D, P, Q, M, and H denote BBB permeability, DRD2 activity, penalized logP, QED, mutagenicity, and human intestinal absorption, respectively. All tasks use the seen-instruction manifests. The IND/OOD distinction follows the paper-declared training scope of property combinations; it does not establish an independent audit of every checkpoint's training examples. Moreover, differing source inclusion filters mean that IND/OOD differences need not be caused solely by unseen combinations.

The backbones are Mistral-7B-Instruct-v0.3 and Llama-3.1-8B-Instruct with their corresponding frozen GeLLMO-P6 adapters. The recorded adapter revisions are \nolinkurl{a6437ea4aedd75584bd9d30001c63e9808489531} (Mistral) and \nolinkurl{6a5b6ab53b8e6d3c2016370a980d57c39c749592} (Llama). The base-model upstream revision fields were not recorded; local model files and source manifests are identified by checksums. The runtime contracts specify bfloat16 base weights, unmerged PEFT adapters with adapter dtype autocasting, and no quantization or offloading. These checkpoints are held fixed across methods within each model.

We retain raw generations, original proposal or trajectory positions, selection provenance, input identifiers, and property-evaluation records. Reusing a cached output requires matching its source input, model, generation protocol, and causal position. A returned output is never chosen by retrospectively maximizing the reported benchmark score unless the method's declared inference rule itself uses that public verifier. The current paired comparison contains nine complete methods and 70,290 task--input outputs. It does not combine different inputs or selectively discard failures to obtain complete rows.

\subsection{Properties, precision, and aggregation}
The same fixed molecular evaluation pipeline supplies the target-property values across methods. BBB permeability, HIA, and mutagenicity use ADMET-AI \citep{admet2024} through the original columns \texttt{BBB\_Martins}, \texttt{HIA\_Hou}, and \texttt{AMES}; DRD2, penalized logP, QED, and synthetic accessibility use the retained GeLLMO property pipeline. These values are computational proxies. They are also public inference tools: methods may query them as specified below. We therefore evaluate complete procedures with different feedback use, rather than claiming that all methods have identical information consumption.

Source-property values remain as stored in the input manifests. Candidate comparison values follow the original pandas 1.0.5/NumPy 1.18.5 pipeline: floating-point \texttt{DataFrame.round(2)}, CSV serialization, and CSV reading. The metric kernel does not round these values again. Directional success requires strict improvement in every target at this comparison precision, with no additional similarity threshold or numerical tolerance. Dataset inclusion thresholds are not success thresholds.

Empty outputs, invalid molecules, and the documented molecule-level failures of the retained evaluation environments count as failures. The accepted exceptions include an initial RDKit parse that succeeds in the inference environment but fails in the legacy evaluation environment, and recorded downstream kekulization failures. Unexpected missing property records, corrupt input identities, and infrastructure failures are unresolved work; they are not silently relabeled as unsuccessful molecules. No failed slot is refilled or replaced by another candidate.

For each task, SR is $100$ times the number of successful returned outputs divided by the full task size. Macro10 averages the ten unrounded task SRs, giving each task equal weight. Because the OOD task sizes differ, this is not the pooled success rate over all 3,905 pairs. Rounded table values are display values. Each main-table row represents one recorded realization of its declared configuration; it does not measure variability over repeated generation seeds.

\subsection{Method contracts}
The methods share the output interface and final evaluator, but retain different internal mechanisms. The following descriptions distinguish the published mechanism from the molecular interface and budget choices used here.

\paragraph{Beam search.}
Beam width and returned generation sequences are both 20, with deterministic decoding and a 100-new-token cap. The benchmark uses the first sequence in the model's beam ordering. It neither searches the beam for a successful molecule nor falls back to a lower-ranked valid candidate. The saved original rank is retained in the endpoint provenance.

\paragraph{Repeated Sampling with verification.}
The generation call draws 40 proposals at temperature 1, top-$p$ 0.95, top-$k$ 50, and a 100-new-token cap. The verification pool is fixed to original positions 0--19, preserving duplicates, invalid strings, and generation order. The verifier is precisely the directional-success predicate, including chemical and property eligibility. The first passing proposal is returned; if none passes, proposal zero is returned unchanged. There is no scalar property score, similarity weight, synthetic-accessibility tie-breaker, or search of positions 20--39. This instantiates the sampling-and-verification mechanism \citep{repeatedsampling2024} with a molecular verifier.

If $y^*$ denotes this output, then $S(x,y^*,T)=\max_{0\leq i<20}S(x,y_i,T)$. Consequently, its $K_{\mathrm{out}}=1$ SR equals the first-20 pool's success coverage by construction. This identity is disclosed rather than interpreted as evidence of a learned selector. The actual generation cost remains 40 proposals; logical verifier queries stop at the first pass and otherwise reach 20. Reading a property cache does not erase the corresponding logical query. Auxiliary output properties, when reported, concern $y^*$ itself.

\paragraph{Self-Refine.}
We preserve the feedback--revision loop \citep{selfrefine2023}, using one refinement at the fixed recorded replicate seed 51001. The returned molecule is the first refined answer, not the best of the initial answer and later rounds. The same model supplies feedback without querying external candidate properties. The reused original prefix contains initial generation, feedback on the initial answer, the first refinement, and terminal feedback on that refinement: four model calls. The terminal feedback does not replace the returned answer. Initial and refinement outputs have a 250-new-token cap; feedback has a 400-token cap. Sampling uses temperature 0.7, top-$p$ 1, and top-$k$ 0. One refinement preserves the method's defining loop, but this configuration does not test the benefit of additional iterations or reproduce a four-refinement budget.

\paragraph{Tree of Thoughts.}
\label{app:tot-scope}
The Tree-of-Thoughts adaptation \citep{tot2023} uses five initial molecular candidates, five language-model voting responses to retain three parents, five new children per selected parent, and five further voting responses to select one of the 15 terminal children. The generator checkpoint also serves as judge. Candidate properties, validity labels, and success labels are absent from the voting prompts; malformed votes contribute zero, and ties use original frontier order. Children have a 160-token cap and judges a 128-token cap, with temperature 1, top-$p$ 1, and top-$k$ 0. This retains causal language-model evaluation of branches and final-frontier readout. The molecular prompts, depth, and width are explicit adaptation choices; neither ToT nor $K_{\mathrm{out}}=1$ requires 20 generated candidates.

This configuration completes terminal-property evaluation on all 3,905 inputs for each model. Tables~\ref{tab:effectiveness}, \ref{tab:workload}, and~\ref{tab:quality} all use this same LM-voted procedure and its actual final outputs.

\paragraph{Graph of Thoughts.}
The G12 molecular adaptation uses the GoT operation framework \citep{got2024}. It generates 12 initial nodes, ranks them, and retains six parents. It then generates six single-parent refinements, six two-parent merges, and four four-parent aggregations, for 28 total nodes. Every second-stage operation uses the selected initial parents and the source molecule. Multi-parent prompts ask the model to combine molecular design ideas into one candidate. They do not perform an explicit chemical graph-edit operation.

The molecular scoring rule is an adaptation defined for this benchmark, not a chemistry formula supplied by the GoT paper. For $m$ target properties, an evaluable node receives
\begin{align}
 s(y)={}&1+\frac{3}{m}\sum_{j=1}^{m}z_j(y)
       +0.50z_{\mathrm{sim}}(y)\nonumber\\
       &+0.25z_{\mathrm{SAS}}(y)-0.25D(y),
\end{align}
where $z_j$ min--max normalizes directional improvement within the current input's candidate pool, $z_{\mathrm{sim}}$ normalizes source similarity, $z_{\mathrm{SAS}}$ rewards lower synthetic-accessibility scores, and $D$ flags a repeated canonical identity. A constant finite feature receives 0.5. Empty and invalid nodes receive total scores $-200$ and $-100$, respectively. Exact score ties preserve creation order. The rule ranks the initial pool for parent selection and independently normalizes all 28 nodes for the final ranking. Parent prompts include the recorded initial-stage feedback. The final output is the original full-graph rank-zero node; no additional common reranker is attached at readout.

Decoding uses temperature 1, top-$p$ 1, top-$k$ 0, with token caps 200 for initial nodes and 220 for later nodes. Similarity is computed from binary Morgan fingerprints with radius 2 and 2,048 bits without chirality. GoT uses raw similarity internally, and the retained comparison precision for target properties and SAS. The topology and score weights are fixed choices of this molecular port; comparisons do not isolate their individual contributions.

\paragraph{s1 budget-forcing component.}
We adapt the inference-time budget-forcing component \citep{s1_2025}, without s1K training or replacement of the GeLLMO checkpoint. Greedy reasoning has a 512-generated-token allowance. At most one literal \texttt{Wait} continuation is inserted when the initial response ends before that allowance is exhausted, using EOS as the local stopping proxy. A separate final-answer cue permits up to 256 additional generated tokens. The parser reads the newly generated final response only; it does not search earlier reasoning for a better molecule. Thus 512 is the reasoning allowance, not the total token cap, and this row is not a reproduction of the full trained s1 system.

\paragraph{Exact Voting.}
Four hundred proposals are generated in 20 calls of 20 sequences each, at temperature 0.1, top-$p$ 1, top-$k$ 0, and a 256-new-token cap. Following the exact-voting mechanism \citep{voting2024}, each proposal casts one vote for its molecular identity. Valid strings are grouped by canonical isomeric SMILES using RDKit 2024.03.6. Invalid responses remain in the vote, grouped by whitespace-normalized response text; repeated invalid responses can share a class, as do empty responses. The largest class wins; tied classes are ordered by first appearance, and the earliest proposal in the winning class is returned. Property quality does not resolve ties. Molecular canonical identity is the task-specific adaptation of exact answer identity.

\paragraph{Four-shot ICL.}
ICL \citep{brown2020} provides a prompting baseline. Four examples are deterministically drawn at seed 2020 from the permitted training pool; each example retains its own task labels. The five held-out property combinations are excluded from that pool according to the paper-declared scope. Both source and target endpoints of demonstrations are excluded if they match any evaluation source molecule, by raw or canonical isomeric SMILES identity. Invalid, empty, duplicate, and unchanged source--target examples are filtered under the recorded demonstration policy. No task-specific similarity retrieval is added. Decoding uses four beams, four returned sequences, length penalty 0.6, and a 100-new-token cap; only the first beam is evaluated. The baseline changes prompt context, rather than implementing a complete published reasoning-time scaling system.

\paragraph{Reflexion.}
The molecular adaptation retains actor generation, evaluator feedback, failure-conditioned verbal reflection, and memory within the same input \citep{reflexion2023}. The evaluator may query the public molecular properties. The process stops at its declared success criterion or after three trials and returns that terminal trial, without retrospectively selecting the best trial. Actor generation uses greedy decoding with a 250-new-token cap. Trial histories are not shared across test inputs. The task-specific evaluator and feedback text are adaptations of the environment interface; the published reflection-and-memory mechanism is preserved.
